%------------------------------------------------------------------------------%
\documentclass[fleqn,utf8,aspectratio=169,12pt]{beamer}

\usepackage{calculo_varias_variaveis-1}

\title{Multiplicadores de Lagrange}

%------------------------------------------------------------------------------%
\begin{document}

\addtitle

%------------------------------------------------------------------------------%
\section{Extremos Condicionados}

%------------------------------------------------------------------------------%
\begin{frame}{Extremos Condicionados}
  Encontre o máximo (mínimo) da função \(f(x, y, z)\)\\[2mm]
  sujeito a \(g(x, y, z) = 0\)

  \pause
  \vspace{\baselineskip}
  \quad \emph{$f$} \quad função objetivo

  \pause
  \quad \emph{$g$} \quad restrição
\end{frame}

%------------------------------------------------------------------------------%
%\framedgraphic{Extremos Condicionados}{campo_gradiente}

%------------------------------------------------------------------------------%
\section{Extremos Locais Restritos a uma Curva}

%------------------------------------------------------------------------------%
\begin{frame}{Teorema do Gradiente Ortogonal}
  Seja \(f(x, y, z)\) uma função diferenciável em uma região

  \pause
  cujo interior contém a curva
  \[
    C\colon\;
    \pause r(t) \;=\; g(t)\veci + h(t)\vecj + k(t)\veck
    \pause \;=\; \vetortri{g(t)}{h(t)}{k(t)}
  \]

  \pause
  Se $P_0$ é um ponto em $C$ onde $f$ possua um máximo ou mínimo local relativo
  aos seus valores em $C$,
  \pause
  então \emph{\(\nabla f\) é normal a $C$ em $P_0$}
\end{frame}

%------------------------------------------------------------------------------%
\begin{proofframe}{Demonstração}
  Os valores de $f$, restrita a $C$, são dados pela função
  \[
    \pause
    p(t) = f(g(t), h(t), k(t))
  \]
  \[
    \pause
    \frac{dp}{dt}
    \pause
    \;=\; \D{f}{g}\frac{dg}{dt}
    + \D{f}{h}\frac{dh}{dt}
    + \D{f}{k}\frac{dk}{dt}
    \pause
    \;=\; \nabla f \cdot \frac{dr}{dt}
  \]

  \pause
  Se $f$, restrita a $C$, possui extremo local em \((a, b, c) = (g(t_0), h(t_0), k(t_0))\)

  \pause
  \vspace{0.4\baselineskip}
  então \(\frac{dp}{dt}(t_0)=0\)
  \pause
  \qquad
  portanto
  \quad
  \emph{\(
    \nabla f(a, b, c) \cdot \frac{dr}{dt}(t_0) = 0
  \)}
\end{proofframe}

%------------------------------------------------------------------------------%
\section{Multiplicadores de Lagrange}

%------------------------------------------------------------------------------%
\begin{frame}{Introdução}
  Para encontrar o máximo (mínimo) da função \(f(x, y, z)\)\\[2mm]
  sujeito a \(g(x, y, z) = 0\)

  \pause
  \vspace{\baselineskip}
  Buscamos pelos pontos onde o gradiente de $f$ \\
  é normal a curva de nível de \(g(x, y, z) = 0\)

  \pause
  \vspace{\baselineskip}
  Isto é, os pontos onde os \emph{gradientes de $f$ e $g$ apontam na mesma direção}
  \pause
  \[
    \nabla f = \lambda\;\nabla g
  \]
  \pause
  \emph{\(\lambda\) é o Multiplicador de Lagrange}
\end{frame}

\framedgraphic{Extremos Condicionados}{LagrangeMultipliers2D}

%------------------------------------------------------------------------------%
\begin{frame}{Método dos Multiplicadores de Lagrange}
  Suponha que \(f(x, y, z)\) e \(g(x, y, z)\) sejam diferenciáveis\\[2mm]
  e \(\nabla g\neq 0\) quando \(g(x, y, z)=0\)

  \pause
  \vspace{\baselineskip}
  Para encontrar os valores extremos locais de \(f(x, y, z)\) \\[2mm]
  \pause
  sujeito a restrição \(g(x, y, z)=0\) \\[2mm]
  \pause
  encontre $x$, $y$, $z$ e $\lambda$ tais que
  \[
    \begin{cases}
      \nabla f = \lambda \nabla g   \\
      g(x, y, z) = 0
    \end{cases}
  \]
\end{frame}

%------------------------------------------------------------------------------%
\section{Exemplos}

\include{J-03-Multiplicadores_Lagrange-ex1}
\include{J-03-Multiplicadores_Lagrange-ex2}

%------------------------------------------------------------------------------%
\section{Funções Crescentes}

%------------------------------------------------------------------------------%
\begin{frame}{Funções Crescentes}
  Se $h(x)$ é uma função real \emph{crescente}
  \pause
  e \(f(x) = h\big(g(x)\big)\) \\
  \pause
  os máximos (mínimos) de $f$ coincidem com os máximos (mínimos) de $g$

  \pause
  \bigskip
  \bigskip
  Se $g$ tem máximo (mínimo) em $a$, $f$ também tem máximo (mínimo) em $a$
  \bigskip
  \bigskip

  \pause
  Note que o valor máximo (mínimo) \emph{não é igual}
\end{frame}

%------------------------------------------------------------------------------%
\section{Exemplos}

\include{J-03-Multiplicadores_Lagrange-ex3}
\include{J-03-Multiplicadores_Lagrange-ex4}
\include{J-03-Multiplicadores_Lagrange-ex5}

%------------------------------------------------------------------------------%
\section{Lista Mínima}

%------------------------------------------------------------------------------%
\begin{frame}{Lista Mínima}
  Cálculo Vol. 2 do Thomas 12\textsuperscript{a} ed. -- Seção 14.8
  \begin{enumerate}
    \item Estudar o texto da seção
    \item Resolver os exercícios: 3, 5, 7, 12, 16, 18, 24 e 30
  \end{enumerate}

  \vfill
  Atenção: A prova é baseada no livro, não nas apresentações
\end{frame}

%------------------------------------------------------------------------------%
\end{document}
%------------------------------------------------------------------------------%
